\section{Arbitrary factor endpoints and stable canonical recovery}
\label{sec:endpoints}

We now use the correspondence classification to recover the multiplication
of an unknown factor.  The argument first realizes that factor, up to
amplification, as a finite-index extension of the given twisted group
factor.  Its self-correspondence grading then consists of bounded
operators and respects the actual product.  This additional multiplication
step is what turns the correspondence data into a group presentation.
The homogeneous-product and minimal-corner construction has a direct
predecessor in Donvil--Vaes \cite[Proposition~6.5 and Theorem~6.6]{DonvilVaes2025}.
Their construction already permits arbitrary source cocycles and recovers
the exact matched ratio. We adapt it from globally defined automorphism
degrees to finite-index isomorphism germs.

\subsection{A finite-index extension and its intrinsic grading}

\begin{lemma}[From a bifinite neighbor to an extension]
\label{lem:endpoints-extension}
Let \(M,P\) be \(\mathrm{II}_1\) factors and let \({}_M K_P\ne0\) be
bifinite.  The commutant \(Q\) of the right \(P\)-action is a positive
finite amplification of \(P\), and the given left action identifies
\(M\) with a unital finite-index subfactor of \(Q\).
\end{lemma}

\begin{proof}
The finite right module has a column realization
\(K=pL^2(P)^n\), with
\(r=(\Tr_n\otimes\tau_P)(p)=\dim K_P>0\).
Its right-action commutant is
\(Q=p(M_n(\mathbb C)\bar\otimes P)p=P^r\), equipped with the
normalized corner trace.  The left action of \(M\) is a normal unital
embedding into \(Q\), because a nonzero representation of a factor has
zero kernel.  Its restriction preserves normalized trace.

As a left \(Q\)-module, \(K\) has finite positive dimension \(1/r\).
Choose an integer \(j\) with \(j/r\geq1\).  Comparison of finite
module projections embeds \(L^2(Q)\) as a left \(Q\)-submodule of
\(K^{\oplus j}\).  Restriction to \(M\) makes it a finite left
\(M\)-module.  Thus the Jones index \([Q:M]=\dim_M L^2(Q)\) is
finite.  This uses the standard finite-factor column and index
identifications; see \cite{Jones1983} and \cite[Theorem~10.1.1, Proposition~10.1.3, and Theorem~10.2.1]{Jones2009}.
\end{proof}

Fix \(\Gamma\in\mathscr K\), an arbitrary normalized scalar cocycle
\(\mu\), and \(M=L_\mu(\Gamma)\).  Write
\(D=\Comm(\Gamma)\) for the group of virtual automorphism germs
defined in Section~\ref{sec:calculus}. The map
\[
 \iota:\Gamma\longrightarrow D,\qquad
 \iota(g)=[\Ad(g)],
\]
is injective: an element in its kernel centralizes a finite-index
subgroup of \(\Gamma\), and hence has a finite conjugacy class.
For a representative \(\delta:A\to B\) of \(d\in D\), one has
\begin{equation}
 d\,\iota(a)\,d^{-1}=\iota(\delta(a))\qquad(a\in A).
 \label{eq:endpoints-inner-germs}
\end{equation}

\begin{lemma}[The intrinsic self-correspondence grading]
\label{lem:endpoints-grading}
Every bifinite \(M\)--\(M\) correspondence has an orthogonal decomposition
\[
 K=\bigoplus_{d\in D}K_d
\]
into finite-dimensional spaces, with only countably many nonzero terms,
such that
\[
 u_gK_du_h^*=K_{\iota(g)d\iota(h)^{-1}}.
\]
The support is contained in finitely many \(\iota(\Gamma)\)-double
cosets.  For every representative \(\delta:A\to B\) of \(d\),
\begin{equation}
 K_d=\left\{\xi\in K:
  \dim_{\mathbb C}\operatorname{span}
    \{u_{\delta(a)}\xi u_a^*:a\in A\}<\infty\right\}.
 \label{eq:input-B1}
\end{equation}
In particular, for the fixed presentation $M=L_\mu(\Gamma)$, the
grading is determined by the two actions and does not depend on a
choice of arithmetic model decomposition.
\end{lemma}

\begin{proof}
Apply Theorem~\ref{thm:correspondences} and the full-germ decomposition
of Section~\ref{sec:calculus}.  Use the inverse of its germ labels, so
that our degree is a germ from the right group to the left group.
If \(x\) is an old label and \(d=x^{-1}\), the old action gives
\[
 \bigl(\Ad(h)x\Ad(g^{-1})\bigr)^{-1}
 =\iota(g)d\iota(h)^{-1}.
\]
Combining the finitely many representation fibres at each identical
inverse label gives \(K_d\) and the stated covariance. The finitely
many old germ orbits become finitely many
\(\iota(\Gamma)\)-double cosets. Separability makes the nonzero
support countable.

For a representative \(\delta\), the operators
\(U_a\xi=u_{\delta(a)}\xi u_a^*\) preserve \(K_d\).
Their multiplier is
\(\mu(\delta(a),\delta(b))\overline{\mu(a,b)}\).
They form a projective unitary representation, so its orbit span is
invariant, even when the multiplier is not trivial.  This proves the
forward inclusion in Equation~\eqref{eq:input-B1}.

Conversely let \(W\) be a finite-dimensional invariant orbit span and
choose an orthonormal basis \(\xi_1,\ldots,\xi_l\) of \(W\).
If \(P_x\) projects onto \(K_x\), the numbers
\[
 m(x)=\sum_{i=1}^l\|P_x\xi_i\|^2
\]
have finite sum \(l\) and are invariant under
\(x\mapsto\iota(\delta(a))x\iota(a)^{-1}\).
Indeed \(U_a\) permutes the degree spaces and carries the chosen basis
of \(W\) to another orthonormal basis.  Positive mass is therefore
possible only on a finite orbit.  The stabilizer of such an \(x\)
contains a finite-index subgroup \(A_0\leq A\), and
\[
 x\iota(a)x^{-1}=\iota(\delta(a))\qquad(a\in A_0).
\]
Represent \(x\) by \(\epsilon:C\to E\) and restrict to
\(A_0\cap C\).  Equation~\eqref{eq:endpoints-inner-germs} and
injectivity of \(\iota\) give \(\epsilon(a)=\delta(a)\) there.
Thus \(x=d\).  All of \(W\) lies in \(K_d\), proving the converse
and the asserted intrinsic description.
\end{proof}

\subsection{Bounded homogeneous operators and the minimal corner}

Let \(M\subset Q\) now be a unital finite-index inclusion of
\(\mathrm{II}_1\) factors and apply Lemma~\ref{lem:endpoints-grading}
to \({}_M L^2(Q)_M\).  Until the next lemma, its degree spaces are only
Hilbert subspaces of \(L^2(Q)\).

\begin{lemma}[Boundedness and multiplication]
\label{lem:endpoints-bounded-grades}
Every \(K_d\) consists of bounded elements of \(Q\), and
\begin{equation}
 K_dK_f\subset K_{df},\qquad K_d^*=K_{d^{-1}},\qquad
 1\in K_1,\qquad u_a\in K_{\iota(a)}.
 \label{eq:endpoints-grade-product}
\end{equation}
Consequently \(K_1\) is a finite-dimensional unital \(*\)-subalgebra
of \(Q\), normalized by the \(u_a\), and \(\tau_Q(K_d)=0\) for
\(d\ne1\).
\end{lemma}

\begin{proof}
Let $\xi\in K_d$. For distinct $g\in\Gamma$, the right translates
$\xi u_g$ have distinct degrees $d\iota(g)$. Injectivity of $\iota$
therefore gives $\langle\xi,\xi u_g\rangle=0$ for $g\ne e$.
Initially put $a=\xi^*\xi\in L^1(Q)_+$. The $L^1$ extension of the
trace expectation satisfies
\[
 E_M(a)=\|\xi\|_2^2 1:
\]
indeed, for each $g\in\Gamma$,
\[
 \tau_M(E_M(a)u_g)=\tau_Q(a u_g)=\langle\xi,\xi u_g\rangle.
\]
These values agree with those of $\|\xi\|_2^2 1$. They determine
the normal functional associated to $E_M(a)\in L^1(M)$, because the
group algebra is ultraweakly dense in $M$. For the bounded spectral
truncations $a_n=\min(a,n)$,
Pimsner--Popa's finite-index inequality
\cite[Proposition~2.1]{PimsnerPopa1986} gives
\[
 [Q:M]^{-1}a_n\leq E_M(a_n)\leq E_M(a)=\|\xi\|_2^2 1.
\]
Monotone convergence makes $a$ bounded, and hence $\xi\in Q$, with
$\|\xi\|_\infty\leq[Q:M]^{1/2}\|\xi\|_2$. This argument uses
orthogonality of the degrees, so the scalar phases in the right action
do not change it. We now verify the product itself.

Take \(\xi\in K_d\), \(\eta\in K_f\), and representatives
\(\delta:A\to B\), \(\epsilon:C\to E\).  On the finite-index
domain \(C_0=\{a\in C:\epsilon(a)\in A\}\), the exact identity
\[
 u_{\delta\epsilon(a)}\xi\eta u_a^*
   =(u_{\delta\epsilon(a)}\xi u_{\epsilon(a)}^*)
      (u_{\epsilon(a)}\eta u_a^*)
\]
puts the orbit of the product into the finite-dimensional span of
products of bases of \(K_d\) and \(K_f\).
These are bounded elements, as just proved.  The intrinsic test
\eqref{eq:input-B1} yields \(\xi\eta\in K_{df}\).
Taking adjoints in the same test gives
\(K_d^*\subset K_{d^{-1}}\), and taking adjoints again gives equality.
The vector \(1\) is fixed by the identity graph.  For the graph of
\(\Ad(a)\), the vector \(u_a\) is carried to a scalar multiple of
itself, since the corresponding product of group labels is \(a\).
Thus its graph orbit span is one-dimensional, even with the twist.
This proves Equation~\eqref{eq:endpoints-grade-product}.

Finite dimension of \(K_1\) comes from the grading lemma, and
normalization follows from its covariance.  Orthogonality to
\(1\in K_1\) gives the trace assertion for other degrees.
\end{proof}

Compressing $K_1$ to a scalar algebra will now force each surviving
degree to be a unitary line.

\begin{lemma}[The minimal identity corner]
\label{lem:endpoints-minimal-corner}
Let \(p\) be a minimal projection of \(K_1\).  There are a countable
subgroup \(\Lambda\leq D\) and a normalized scalar cocycle \(\omega\)
such that the actual corner \(pQp\), with normalized trace, is
isomorphic to \(L_\omega(\Lambda)\).  Each nonzero space \(pK_dp\)
is one-dimensional and has a unitary basis vector in the corner.
\end{lemma}

\begin{proof}
Left and right multiplication by \(p\) are commuting orthogonal
projections on \(L^2(Q)\).  They preserve each degree by
Equation~\eqref{eq:endpoints-grade-product}; their product projects
onto \(L^2(pQp)\).  Thus, up to the common normalization factor of
the corner trace,
\[
 L^2(pQp)=\bigoplus_{d\in D}pK_dp,\qquad pK_1p=\mathbb Cp.
\]
For \(0\ne x\in pK_dp\), the product and adjoint laws give
\(x^*x=\alpha p\) and \(xx^*=\beta p\), with positive scalars.
Traciality gives \(\alpha=\beta\).  Therefore
\(x/\sqrt\alpha\) is a unitary in \(pQp\).
If \(y\) belongs to the same degree, \(x^*y\in\mathbb Cp\), so
\(y\) is a scalar multiple of \(x\).

Set \(\Lambda=\{d:pK_dp\ne0\}\).  It contains the identity, is
closed under inverses, and is closed under products because products
of its homogeneous unitaries cannot vanish.  It is countable by
orthogonality and separability.  Choose unitary vectors
\(v_d\in pK_dp\), with \(v_1=p\), and define
\[
 v_dv_f=\omega(d,f)v_{df}.
\]
Unitarity makes the coefficient unimodular, associativity gives the
2-cocycle identity, and \(v_1=p\) gives normalization.
The \(v_d\) form an orthonormal basis for the normalized corner
trace.  The unitary \(\delta_d\mapsto v_d\) from \(\ell^2(\Lambda)\)
to \(L^2(pQp)\) intertwines twisted left regular multiplication with
left multiplication by these units.

Let \(Q_0\) be the von Neumann algebra they generate.  The
trace-preserving expectation onto \(Q_0\) has \(L^2\)-range containing
every \(v_d\), so it has full range and \(Q_0=pQp\).
The basis unitary therefore gives a normal unital complex-linear
\(*\)-isomorphism \(L_\omega(\Lambda)\cong pQp\), as claimed.
\end{proof}

\subsection{The ratio condition and the reconstructed group}

The minimal corner has supplied a twisted group algebra, but the group
$\Lambda$ still needs to satisfy the hypotheses of
Theorem~\ref{thm:endpoints}. We next recover a finite-index copy of a
subgroup of $\Gamma$ inside $\Lambda$, obtain the projective ratio from
the identity matrix block, and prove that $\Lambda$ is ICC.

\begin{proposition}[Arithmetic constraints on the corner]
\label{prop:endpoints-corner-constraints}
For the corner in Lemma~\ref{lem:endpoints-minimal-corner}, the group
\(\Lambda\) is countably infinite and ICC.  There are finite-index
subgroups \(A\leq\Gamma\), \(B\leq\Lambda\), an actual group
isomorphism \(\delta:A\to B\), and a finite-dimensional projective
representation with multiplier
\(\mu|_A/\delta^*(\omega|_B)\).
\end{proposition}

\begin{proof}
Let \(z\) be the central support of \(p\) in the finite-dimensional
algebra \(K_1\), so that \(zK_1\) is the matrix block containing
\(p\). We first split the action on this block from the homogeneous
unitaries that will survive in \(pQp\).
The subgroup
\[
 A=\{a\in\Gamma:u_azu_a^*=z\}
\]
has finite index, because conjugation permutes the finitely many
central blocks. Although these conjugations preserve $z$, they may move
$p$ inside its matrix block. We therefore remove their action on that
block before compressing. Choose unitary matrix implementers
\(s_a\in\mathcal U(zK_1)\), \(s_e=z\), for
\(\Ad(u_a)|_{zK_1}\).  Since these automorphisms form an ordinary
group action, there is a normalized scalar cocycle \(\kappa\) with
\[
 s_as_b=\kappa(a,b)s_{ab}.
\]
These matrices already constitute a finite-dimensional projective
representation.  Their image is allowed to be infinite.

The unitaries \(q_a=s_a^*u_az\in zQz\) commute with \(zK_1\).
Using \(u_as_b^*u_a^*=s_as_b^*s_a^*\), their product is
\begin{align*}
 q_aq_b
  &=s_a^*(u_as_b^*u_a^*)u_au_bz\\
  &=\mu(a,b)s_b^*s_a^*u_{ab}z
    =\mu(a,b)\overline{\kappa(a,b)}q_{ab}.
\end{align*}
They are homogeneous of degree \(\iota(a)\), and they commute with
\(p\).  Their compressions \(t_a=pq_ap\) are therefore corner
unitaries in \(pK_{\iota(a)}p\), satisfying the same product formula.
It follows that \(B=\iota(A)\) lies in \(\Lambda\), and that
\(\delta=\iota|_A:A\to B\) is an actual isomorphism.

To compare their multiplier with the cocycle already chosen in
Lemma~\ref{lem:endpoints-minimal-corner}, write the basis unitary of
degree \(\iota(a)\) as
\(v_{\iota(a)}=\lambda(a)t_a\), where $\lambda:A\to\mathbb T$
and $\lambda(e)=1$. With the convention
\[
 (\partial\lambda)(a,b)=\lambda(a)\lambda(b)\lambda(ab)^{-1},
\]
we obtain
\begin{equation}
 \delta^*\omega=(\partial\lambda)\mu|_A\,\overline\kappa,
 \qquad
 \frac{\mu|_A}{\delta^*\omega}
       =(\partial\lambda)^{-1}\kappa.
 \label{eq:endpoints-ratio}
\end{equation}
On the defining space of the block, the matrices
\(a\mapsto\overline{\lambda(a)}s_a\) realize precisely the
second multiplier.  This is the required ratio condition.

We have the subgroup isomorphism and its exact ratio. It remains to
show that its image has finite index in \(\Lambda\) and that
\(\Lambda\) is ICC. Put \(H=\iota(\Gamma)\).
Every \(d\in D\) commensurates \(H\), by
Equation~\eqref{eq:endpoints-inner-germs} on a finite-index domain.
Consequently each double coset \(HdH\) is a finite union of right
\(H\)-cosets.  Lemma~\ref{lem:endpoints-grading} puts the original
degree support in finitely many such double cosets.  Since \(B\) has
finite index in \(H\), that support is contained in finitely many
right \(B\)-cosets.  The subgroup \(\Lambda\) is contained in this
support and contains \(B\); hence \([\Lambda:B]<\infty\).
In particular it is infinite.

Suppose \(l\in\Lambda\) has finite conjugacy class. Its centralizer
$C_\Lambda(l)$ has finite index in \(\Lambda\). Thus
$A_0=\iota^{-1}(C_\Lambda(l)\cap B)$ has finite index in
$\Gamma$, and $l$ centralizes \(\iota(A_0)\).
Represent its germ by \(\epsilon:C\to E\).  For
\(a\in A_0\cap C\), centralization and
Equation~\eqref{eq:endpoints-inner-germs} give
\(\iota(\epsilon(a))=\iota(a)\).
Injectivity of \(\iota\) gives \(\epsilon(a)=a\) on this
finite-index subgroup.  Thus \(l=1\) in \(D\), proving ICC.
The construction uses only subgroup isomorphisms and introduces no
finite kernel.
\end{proof}

\begin{proof}[Proof of Theorem~\ref{thm:endpoints}]
Apply Lemma~\ref{lem:endpoints-extension} to the given nonzero bifinite
correspondence. It gives a finite-index extension
\(M\subset Q=P^r\).  The preceding lemmas and
Proposition~\ref{prop:endpoints-corner-constraints} give
\(pQp\cong L_\omega(\Lambda)\) with the asserted group and ratio
data.  Since \(pQp\) is the amplification
\(P^{r\tau_Q(p)}\), standard column equivalences give
\[
 P\cong L_\omega(\Lambda)^t,
 \qquad t=(r\tau_Q(p))^{-1}>0.
\]

Conversely the stated subgroup and projective data give a nonzero
bifinite arithmetic correspondence from \(M\) to
\(L_\omega(\Lambda)\), by Section~\ref{sec:calculus}.
Compose it with the corner equivalence from that factor to its
amplification and then with the specified factor isomorphism to \(P\).
This is a nonzero bifinite \(M\)--\(P\) correspondence.

When \(\mu=1\), conjugating the finite-dimensional ratio
representation and transporting it through \(\delta\) gives a
representation with multiplier \(\omega|_B\).  Projective induction
from \(B\) to \(\Lambda\), with the given global multiplier
\(\omega\), is finite-dimensional and realizes \(\omega\) itself;
compare \cite[Lemmas~6.7--6.8]{DonvilVaes2025}. The induction formula of Section~\ref{sec:calculus} applies equally to
this finite-index subgroup.  For general \(\mu\) the construction
proves only the exact ratio in Equation~\eqref{eq:endpoints-ratio},
as required.
\end{proof}

\subsection{Stable scales and the actual isomorphism}

For a specified stable isomorphism, the category gives a different
constraint: its equivalence correspondence has reciprocal dimensions.
After integrality forces both dimensions to be one, we will use the
given corner columns to identify the specified map. All dimensions below
use the normalized factor traces and the ordinary unnormalized matrix
trace $\Tr_n$, with $\Tr_n(1)=n$.

\begin{lemma}[Dimension and corner equivalences]
\label{lem:endpoints-equivalence-dimensions}
For composable bifinite correspondences of finite factors, left
dimensions multiply under relative tensor product, and so do right
dimensions.  If
\(M^t=p(M_n(\mathbb C)\bar\otimes M)p\) and
\(C_p=pL^2(M)^n\), then \({}_{M^t}C_p{}_M\) is an equivalence
correspondence with dimensions \((1/t,t)\) in left/right order.
\end{lemma}

\begin{proof}
For the right-dimension assertion, realize a right \(N\)-module
\(H\) as \(eL^2(N)^k\).  Then
\(H\otimes_N K=eK^k\).  The normalized trace on the finite factor
\(\operatorname{End}(K_P)\) restricts to the normalized trace on its
unital copy of \(N\).  Computing the dimension of \(eK^k\) therefore
gives
\[
 \dim(H\otimes_N K)_P
  =(\Tr_k\otimes\tau_N)(e)\,\dim K_P.
\]
The left formula follows by conjugating the correspondences.
The column \(C_p\) has right dimension \(t\) by its defining
projection.  The standard column pairings identify
\(C_p\otimes_M\overline{C_p}\) with \(L^2(M^t)\) and
\(\overline{C_p}\otimes_{M^t}C_p\) with \(L^2(M)\).
The multiplication formula makes its left dimension \(1/t\).
These identifications use the normalized corner trace, which accounts
for the reciprocal scale.
\end{proof}

\begin{proof}[Proof of Theorem~\ref{thm:stable}]
Set \(M=L_\omega(\Lambda)\) and \(N=L_\nu(\Pi)\).
Use the actual given corner presentations
\[
 M^t=pM_n(M)p,\qquad N^s=qM_m(N)q,
\]
and the columns \(C_p=pL^2(M)^n\), \(C_q=qL^2(N)^m\).
A specified isomorphism \(\theta:M^t\to N^s\) gives the
\(M\)--\(N\) equivalence correspondence
\begin{equation}
 H_\theta=\overline{C_p}\otimes_{M^t}
 {}_\theta L^2(N^s)\otimes_{N^s}C_q.
 \label{eq:endpoints-lifted-isomorphism}
\end{equation}
By Lemma~\ref{lem:endpoints-equivalence-dimensions}, its left and right
dimensions are \(t/s\) and \(s/t\). Theorem~\ref{thm:correspondences}
applies to \(H_\theta\) because \(\Lambda\in\mathscr K\) and
\(\Pi\) is countably infinite ICC. Its full-germ decomposition and
Equation~\eqref{eq:calculus-summand-dimensions} make both numbers
positive integers. Their product is one, so \(t=s\).

The same formula now forces a single full-germ orbit, a
one-dimensional fibre, and both stabilizer projections equal to the
whole endpoint groups.  Indeed each nonzero orbit contributes its
positive integer fibre dimension times a positive integer subgroup
index to each of the two sums, both of which equal one.  Thus its
graph is an actual group isomorphism \(\delta:\Lambda\to\Pi\).
Consequently \(\Pi\in\mathscr K\), since \(\Lambda\in\mathscr K\)
and this class is invariant under group isomorphism.
The one-dimensional fibre is a normalized cochain \(b:\Lambda\to
\mathbb T\) with
\[
 \omega(g,h)=b(g)b(h)b(gh)^{-1}\nu(\delta(g),\delta(h)).
\]
The arithmetic bundle identification in Section~\ref{sec:calculus}
therefore gives
\(H_\theta\cong {}_\varphi L^2(N)\), where
\[
 \varphi(u_g)=b(g)v_{\delta(g)}.
\]
Multiplication and the displayed cochain equation define a
\(*\)-homomorphism on the group algebra.  Its trace-preserving basis
unitary extends it to a normal unital isomorphism \(M\to N\).

We now return to the specified map in its actual corners. First choose
any partial isometry between the corresponding projections; then the
correspondence identification will supply the inner correction that makes
it implement \(\theta\). The projections
\(\varphi^{(n)}(p)\) and \(q\) have equal unnormalized trace \(t=s\).
Comparison of projections in the finite matrix amplification of
\(N\) gives a rectangular partial isometry
\[
 z\in M_{m,n}(N),\qquad
 z^*z=\varphi^{(n)}(p),\qquad zz^*=q.
\]
Define
\(\theta_0(x)=z\varphi^{(n)}(x)z^*\) on \(pM_n(M)p\).
Its lifted correspondence in
Equation~\eqref{eq:endpoints-lifted-isomorphism} is
\({}_\varphi L^2(N)\).  To check this directly, on bounded columns
set $F(\xi)=z\varphi^{(n)}(\xi)$, applying $\varphi$ entrywise.
For $x\in pM_n(M)p$, $a\in M$, and $\xi,\eta\in pM^{n}$, one has
\[
 F(x\xi a)=\theta_0(x)F(\xi)\varphi(a),\qquad
 F(\xi)^*F(\eta)=\varphi(\xi^*\eta).
\]
Moreover $F$ is onto the bounded columns $qN^m$, since $zz^*=q$.
These identities, followed by the normalized column pairings of
Lemma~\ref{lem:endpoints-equivalence-dimensions}, contract the two
column factors to the trace module with left action \(\varphi\).

Consequently \(H_\theta\cong H_{\theta_0}\).  Tensor on the left
with \(C_p\) and on the right with \(\overline{C_q}\); the column
equivalences cancel and give a unitary
\[
 U:{}_\theta L^2(N^s)\longrightarrow {}_{\theta_0}L^2(N^s).
\]
A right-linear unitary on this standard trace module is left
multiplication by a unitary \(w\in N^s\).
Its left intertwining identity says
\(w\theta(x)=\theta_0(x)w\).
Replacing \(z\) by \(w^*z\), which has the same initial and final
projections, therefore yields the actual formula
\begin{equation}
 \theta(x)=z\varphi^{(n)}(x)z^*
 \qquad(x\in pM_n(M)p).
 \label{eq:endpoints-actual-corner-map}
\end{equation}
At ordinary unit scale this is precisely
\(\theta(u_g)=z\,b(g)v_{\delta(g)}z^*\) with \(z\in N\) unitary.

Conversely a group isomorphism and normalized cochain satisfying the
cocycle equation construct \(\varphi\) by the same basis unitary.
At equal scales, comparison of the two actual projections supplies
\(z\), and Equation~\eqref{eq:endpoints-actual-corner-map} constructs
the corner isomorphism.  This proves both directions and the
specified-map assertion.  Taking identical factor inputs and scales
\(1,t\) shows that every factor in the statement has fundamental
group \(\{1\}\).
\end{proof}
